%% Abstract and introduction text from Phase I.
This paper presents the design, implementation and functional verification of a
32-bit, five-stage pipelined processor for a subset of the RISC-V RV32I base integer
instruction set. The processor is written in synthesizable-style Verilog HDL and is
organised into the classic Instruction Fetch (IF), Instruction Decode (ID), Execute (EX),
Memory (MEM) and Write-Back (WB) stages, separated by four pipeline registers. It supports
29 instructions covering register and immediate arithmetic, logical and shift operations,
word loads and stores, conditional branches, jumps and upper-immediate loads. Data hazards
are resolved by a forwarding unit that bypasses results from the EX/MEM and MEM/WB
registers to the ALU, a hazard detection unit that inserts a single-cycle bubble for
load-use dependencies, and a write-first register file. Control hazards are handled by
resolving branches and jumps in the EX stage under a predict-not-taken policy, flushing
two wrongly fetched instructions on a taken transfer. Structural hazards are avoided by
separate instruction and data memories. The design was verified with Icarus Verilog using
16 self-checking unit testbenches (204 checks) and six system-level assembly programs
(36 architectural-state checks), all of which pass. A testbench performance counter
measures cycles per instruction (CPI) on every program; including pipeline fill, the
measured CPI ranges from 1.19 to 2.09, and every cycle count matches the analytical
model $N + 4 + S + 2T$ for $N$ instructions, $S$ load-use stalls and $T$ taken branches
or jumps.

\section{Introduction}
A single-cycle processor completes every instruction in one long clock cycle, so the clock
period is set by the slowest instruction. Pipelining divides instruction execution into
stages and overlaps the stages of consecutive instructions, in the same way as an assembly
line. In the ideal case one instruction completes every cycle, and the clock period is set
by the slowest \emph{stage} rather than the slowest instruction~\cite{pattersonRV,hennessyQA}.

Overlapping instructions introduces \emph{hazards}: an instruction may need a result that
has not yet been written back (data hazard), the next instruction address may not yet be
known (control hazard), or two stages may need the same resource (structural hazard).
A correct pipelined processor must detect and resolve all three.

RISC-V is an open, royalty-free instruction set architecture developed at the University
of California, Berkeley~\cite{rvmanual,asanovic2014free}. Its regular, fixed-length
32-bit encoding and small base integer set make it well suited to teaching and to
hardware design projects. This project implements an RV32I subset on a five-stage
pipeline with full forwarding and hazard handling, and verifies it in simulation.

The main contributions of this work are:
\begin{itemize}
  \item a modular Verilog implementation of a five-stage RV32I pipeline whose modules
        share one header of opcodes and control codes;
  \item forwarding, load-use stall and branch-flush logic that together resolve all
        data and control hazards for the supported instructions;
  \item a self-checking verification suite of unit testbenches and assembly programs,
        run by a single script.
\end{itemize}

Section~II reviews related work. Section~III describes the proposed design.
Section~IV reports verification and performance results, and Section~V lists
individual contributions.

% ---------------------------------------------------------------------------
